%%%%%%%%%%% En-tête %%%%%%%%%%%%%
\documentclass[11pt]{article}		            % Précise le type de document, et la taille de la police de caractère
\usepackage[margin=2.5cm]{geometry}             % Précise les marges du document
\usepackage{amsmath,amsthm,amssymb,amsfonts}	% Pour pouvoir inclure certains symboles et environnements mathématiques
\usepackage[french]{babel}	                    % Pour préciser la langue du document
\usepackage[T1]{fontenc}                        % accents codés dans la fonte
\usepackage[utf8]{inputenc}                        % accents codés dans la fonte
\usepackage{graphicx}	                        % Pour inclure des images
\usepackage{color}	                            % Pour inclure du texte en couleur
\usepackage[colorlinks]{hyperref}                           % Pour inclure des liens (URL)
\newcommand{\N}{\mathbb{N}}	                    % Commande personnelle, plus rapide pour taper les ensembles
\newcommand{\Z}{\mathbb{Z}}	                    % Commande personnelle, plus rapide pour taper les ensembles
\newcommand{\R}{\mathbb{R}}	                    % Commande personnelle, plus rapide pour taper les ensembles
\theoremstyle{definition}                       % Définition d'environnements
    \newtheorem{defn}{Définition}
    \newtheorem{thm}{Théorème}
    \newtheorem{lemme}{Lemme}
    \newtheorem*{dem*}{Dém}
    \newtheorem{rem}{Remarque}
    \newtheorem{prop}{Proposition}
    \newtheorem{ex}{Exemple}
    \newtheorem{question}{Question}
%%%%%%%%%%%%%%%%%%%%%%%%%%%%%%%%%

\begin{document}                                % Début du document
\thispagestyle{empty}	                        % Pour éviter d'avoir un en-tête et un pied de page sur la page couverture
\begin{center}
    \includegraphics[width=5cm]{logo.png}	    % Pour inclure le logo (on précise la largeur de l'image)
\end{center}
\vspace{4cm}	                                % Espacement vertical
\begin{center}	                                % On centre le texte
    {\huge \bf Titre }\\                        % \huge fait que le texte est gros, \bf fait que le texte est gras
    \vspace{3cm}
    \large Travail présenté à M. Untel dans le cadre du cours \emph{XXX-YYYY : Nom du cours} \\
    \vspace{3cm}
    Travail réalisé par A et B\\
    \vfill	                                    % On va jusqu'au bas de la page avant de mettre le texte ci-dessous
    date de remise
\end{center}
\newpage
\begin{question}
    Ceci est la première question. Ceci est la première question. Ceci est la première question. Ceci est la première question. Ceci est la première question.\\
    
    Pour y répondre, je dois écrire des équations :
    \begin{align*}
        a&=b\\
        c&=\sqrt{x}
    \end{align*}
\end{question}
\begin{question}
    Dans cette question, nous devons trouver une photo du père des processus stochastiques. La voici:
    \begin{figure}[!h]      % Le !h entre crochet c'est pour forcer la position de la photo ici.
        \centering
        \includegraphics[width=10cm]{AAMarkov.jpg}
        \caption{Andreï Andreïevitch Markov (1856-1922) (voir \cite{wiki:Markov})}
        \label{fig:markov}
    \end{figure}
\end{question}
Voici un tableau.
\begin{table}[h!]
    \centering
    \begin{tabular}{|c|c|c|c|}
        \hline
          & 1    & 2    & 3 \\ \hline
        1 & 1    & 2    & 3 \\ \hline
        2 & 2    & 4    & 6 \\ \hline
    \end{tabular}
\end{table}

\newpage
    Voici une matrice identité:
    \[
        \begin{pmatrix}
         1    &0    & 0 \\
         0    & 1    & 0 \\
         0    & 0    & 1 \\
        \end{pmatrix}
    \]
    \begin{center}
        Voici du texte au centre.
    \end{center}
    \begin{flushright}
                {\color{red}Voici du texte à droite en rouge.}
    \end{flushright}
Pour en apprendre sur la théorie des jeux, vous pouvez consulter l'ouvrage de James Webb (voir \cite{Webb}).
\newpage
\bibliography{bibliographie.bib} 
\bibliographystyle{unsrt}
\end{document}